% Generated from locale/tr/content/first-order-logic/completeness/construction-of-model.tex; do not hand-edit.


\iftag{FOL}
      {\olfileid[tr]{fol}{com}{mod}}
      {\olfileid[tr]{pl}{com}{mod}}

\olsection{Bir Modelin Kuruluşu}

\begin{explain}
\iftag{FOL}{Şimdilik~$\eq$ ile ilgilenmiyoruz; yani yalnızca~$\eq$
  içermeyen tutarlı bir !!{sentence}s kümesi~$\Gamma$'nın
  sağlanabilir olduğunu göstermek istiyoruz. Önce~$\Gamma$'yı tutarlı,
  !!{complete} ve doygun bir~$\Gamma^*$ kümesine genişletiriz. Bu
  durumda $\Struct{M(\Gamma^*)}$ modelinin tanımı basittir:
  $\Lang{L'}$ dilinin kapalı terimleri kümesini evren olarak alırız.
  Her !!{constant} simgesini kendisine atar ve daha genel olarak her
  kapalı~$t$ terimi için $\Value{t}{M(\Gamma^*)}=t$ olmasını sağlarız.
  !!{predicate}s simgelerine, atomik bir !!{sentence}
  $\Struct{M(\Gamma^*)}$ içinde ancak ve ancak~$\Gamma^*$ içindeyse
  doğru olacak biçimde uzanımlar atanır. Bu,~$\Gamma^*$ içindeki bütün
  atomik !!{sentence}s ifadelerini açıkça $\Struct{M(\Gamma^*)}$
  içinde doğru kılar. Geri kalanlar da başlangıçtaki~$\Gamma^*$
  tutarlı, tam ve doygunsa doğrudur.}{Şimdi bütün $!A\in\Gamma$
  ifadelerini doğru kılan !!a{valuation} tanımlamaya hazırız.
  Bunun için önce Lindenbaum Lemması'nı uygularız ve
  $\Gamma^*\supseteq\Gamma$ olacak biçimde tam ve tutarlı bir
  $\Gamma^*$ elde ederiz. $\Gamma^*$ içindeki !!{propositional variable}s
  ifadelerinin~$\pAssign v(\Gamma^*)$ değerlemesini
  belirlemesine izin veririz.}
\end{explain}

\iftag{FOL}{%
\begin{defn}[Terim modeli]
\ollabel{defn:termmodel}
$\Gamma^*$, bir~$\Lang L$ dilindeki !!a{complete}, tutarlı ve doygun
bir !!{sentence}s kümesi olsun. $\Gamma^*$ kümesinin \emph{terim
  modeli}~$\Struct M(\Gamma^*)$ aşağıdaki gibi tanımlanan
!!{structure} ifadesidir:
\begin{enumerate}
\item !!{domain}~$\Domain{M(\Gamma^*)}$,~$\Lang L$ dilinin bütün kapalı
  terimleri kümesidir.
\item !!a{constant}~$c$ simgesinin yorumu kendisidir:
  $\Assign{c}{M(\Gamma^*)}=c$.
\item !!{function}~$f$ simgesine, argüman olarak kapalı $t_1$, \dots,
  $t_n$ terimleri verildiğinde değer olarak kapalı
  $f(t_1,\dots,t_n)$ terimini veren fonksiyon atanır:
\[
\Assign{f}{M(\Gamma^*)}(t_1, \dots, t_n) = f(t_1,\dots, t_n)
\]
\item $R$, $n$-li !!{predicate} ise
  \[
  \tuple{t_1, \dots,t_n}\in\Assign{R}{M(\Gamma^*)}
  \text{ ancak ve ancak }
  \Atom{R}{t_1,\dots,t_n}\in\Gamma^*.
  \]
\end{enumerate}
\end{defn}
}{%
\begin{defn}
\ollabel{defn:termmodel}
$\Gamma^*$'ın tam ve tutarlı bir !!{formula}s kümesi olduğunu
varsayalım. Şöyle tanımlarız:
\[
\pAssign v(\Gamma^*)(p) = \begin{cases}
  \True & \text{$p\in\Gamma^*$ ise}\\
  \False & \text{$p\notin\Gamma^*$ ise.}
  \end{cases}
\]
\end{defn}
}

\iftag{FOL}{%
Şimdi gerçekten $\Value{t}{M(\Gamma^*)}=t$ olduğunu doğrulayalım.

\begin{lem}
\ollabel{lem:val-in-termmodel} $\Struct M(\Gamma^*)$,
\olref{defn:termmodel} tanımındaki terim modeli olsun. O zaman her kapalı
$t$ terimi için $\Value{t}{M(\Gamma^*)}=t$.
\end{lem}
\begin{proof}
İspat~$t$ üzerinde tümevarımladır. $t$ bir !!a{constant} olduğundaki
temel durum doğrudan terim modeli tanımından çıkar. Tümevarım adımında
$t_1,\ldots,t_n$ kapalı terimlerinin
$\Value{t_i}{M(\Gamma^*)}=t_i$ eşitliklerini sağladığını ve $f$'nin
$n$-li bir !!{function} olduğunu varsayalım. O zaman
\begin{align*}
\Value{f(t_1,\ldots,t_n)}{M(\Gamma^*)} &= \Assign{f}{M(\Gamma^*)}(\Value{t_1}
 {M(\Gamma^*)},
\ldots, \Value{t_n}{M(\Gamma^*)}) \\
&= \Assign{f}{M(\Gamma^*)}(t_1, \dots, t_n) \\
&= f(t_1,\dots, t_n).
\end{align*}
Dolayısıyla tümevarım gereği eşitlik her kapalı~$t$ terimi için
geçerlidir.
\end{proof}
}

\iftag{FOL}{%
\begin{explain}
  \iftag{prvEx}{!!^a{structure}~$\Struct{M}$, hiçbir doğru
    $!A(t)$ örneği bulunmadığı halde varlıksal niceleyicili bir
    $\lexists[x][!A(x)]$ !!{sentence} ifadesini doğru kılabilir.}{}
  \iftag{prvAll}{!!^a{structure}~$\Struct{M}$, tümel niceleyicili bir
    $\lforall[x][!A(x)]$ !!{sentence} ifadesinin bütün $!A(t)$
    örneklerini doğru kıldığı halde $\lforall[x][!A(x)]$ ifadesini
    doğru kılmayabilir.}{} Bunun nedeni, genel olarak
  $\Domain{M}$ içindeki her !!{element} öğesinin kapalı bir terimin
  değeri olmamasıdır ($\Struct{M}$ örtülmüş olmayabilir). Sağlama
  bağıntısının değişken atamaları üzerinden tanımlanmasının nedeni
  budur. Ancak terim modelimiz~$\Struct{M(\Gamma^*)}$ için buna gerek
  kalmaz; çünkü bu model örtülmüştür. Bir sonraki sonuç tam olarak bunu
  söyler.
\end{explain}

\begin{prop}
\ollabel{prop:quant-termmodel}
$\Struct M(\Gamma^*)$,~\olref{defn:termmodel} tanımındaki terim modeli
olsun.
\begin{tagenumerate}{prvEx,prvAll}
\tagitem{prvEx}{$\Sat{M(\Gamma^*)}{\lexists[x][!A(x)]}$ ancak ve ancak
  en az bir kapalı~$t$ terimi için $\Sat{M(\Gamma^*)}{!A(t)}$ ise.}{}
\tagitem{prvAll}{$\Sat{M(\Gamma^*)}{\lforall[x][!A(x)]}$ ancak ve ancak
  bütün kapalı~$t$ terimleri için $\Sat{M(\Gamma^*)}{!A(t)}$ ise.}{}
\end{tagenumerate}
\end{prop}

\begin{proof}
\begin{tagenumerate}{prvEx,prvAll}
\tagitem{prvEx}{%
  \iftag{probEx}{Alıştırma.}{\olref[syn][ass]{prop:sat-quant} uyarınca
    $\Sat{M(\Gamma^*)}{\lexists[x][!A(x)]}$ ancak ve ancak en az bir
    değişken ataması~$s$ için $\Sat{M(\Gamma^*)}{!A(x)}[s]$ ise.
    $\Domain{M(\Gamma^*)}$,~$\Lang{L}$ dilinin kapalı terimlerinden
    oluştuğundan bu, ancak ve ancak $s(x)=t$ ve
    $\Sat{M(\Gamma^*)}{!A(x)}[s]$ olacak biçimde en az bir kapalı~$t$
    terimi varsa geçerlidir. \olref[fol][syn][ext]{prop:ext-formulas}
    uyarınca, $s(x)=t$ olduğunda $\Sat{M(\Gamma^*)}{!A(x)}[s]$ ancak
    ve ancak $\Sat{M(\Gamma^*)}{!A(t)}[s]$ ise.
    \olref[fol][syn][ass]{prop:sentence-sat-true} uyarınca $!A(t)$ bir
    tümce olduğundan $\Sat{M(\Gamma^*)}{!A(t)}[s]$ ancak ve ancak
    $\Sat{M(\Gamma^*)}{!A(t)}$ ise.}}{}
\tagitem{prvAll}{%
  \iftag{probAll}{Alıştırma.}{\olref[syn][ass]{prop:sat-quant} uyarınca
    $\Sat{M(\Gamma^*)}{\lforall[x][!A(x)]}$ ancak ve ancak her değişken
    ataması~$s$ için $\Sat{M(\Gamma^*)}{!A(x)}[s]$ ise.
    $\Domain{M(\Gamma^*)}$ kümesinin~$\Lang{L}$ dilinin kapalı
    terimlerinden oluştuğunu hatırlayalım. Bu yüzden her kapalı~$t$
    terimi için $s(x)=t$ olacak bir değişken ataması vardır ve her
    değişken atamasında $s(x)$, kapalı bir~$t$ terimidir.
    \olref[fol][syn][ext]{prop:ext-formulas} uyarınca, $s(x)=t$
    olduğunda $\Sat{M(\Gamma^*)}{!A(x)}[s]$ ancak ve ancak
    $\Sat{M(\Gamma^*)}{!A(t)}[s]$ ise.
    \olref[fol][syn][ass]{prop:sentence-sat-true} uyarınca $!A(t)$ bir
    tümce olduğundan $\Sat{M(\Gamma^*)}{!A(t)}[s]$ ancak ve ancak
    $\Sat{M(\Gamma^*)}{!A(t)}$ ise.}}{}
\end{tagenumerate}
\end{proof}
}{}

\tagprob[FOL]{probEx,probAll}
\begin{prob}
\olref[fol][com][mod]{prop:quant-termmodel} önermesinin ispatını tamamlayın.
\end{prob}
\tagendprob

\begin{lem}[Doğruluk Lemması]
\ollabel{lem:truth} \iftag{FOL}{$!A$ ifadesinin~$\eq$ içermediğini
varsayalım. O zaman}{}
$\iftag{FOL}{\Sat{M(\Gamma^*)}}{\pSat{v(\Gamma^*)}}{!A}$ ancak ve ancak
$!A\in\Gamma^*$ ise.
\end{lem}

\begin{proof}
İki yönü aynı anda,~$!A$ üzerinde tümevarımla ispatlarız.
\begin{enumerate}
\tagitem{prvFalse}{\indcase{!A}{\lfalse}
  {$\iftag{FOL}{\Sat/{M(\Gamma^*)}{\lfalse}}{\pSat/{v(\Gamma^*)}{\lfalse}}$
    sağlama tanımı gereği geçerlidir. Öte yandan $\Gamma^*$ tutarlı
    olduğundan $\lfalse\notin\Gamma^*$.}}{}

\tagitem{prvTrue}{\indcase{!A}{\ltrue}
  {$\iftag{FOL}{\Sat{M(\Gamma^*)}}{\pSat{v(\Gamma^*)}}{\ltrue}$
    sağlama tanımı gereği geçerlidir. Öte yandan $\Gamma^*$ tutarlı ve
    !!{complete} olduğundan ve $\Gamma^*\Proves\ltrue$ olduğundan
    $\ltrue\in\Gamma^*$.}}{}

\tagitem{FOL}{\indcase{!A}{R(t_1, \dots, t_n)}
  {$\Sat{M(\Gamma^*)}{\Atom{R}{t_1, \dots, t_n}}$ ancak ve ancak
    $\tuple{t_1,\dots,t_n}\in\Assign{R}{M(\Gamma^*)}$ ise (sağlama
    tanımı gereği); bu da ancak ve ancak $R(t_1,\dots,t_n)\in\Gamma^*$
    ise geçerlidir ($\Struct M(\Gamma^*)$ kuruluşu gereği).}}{\indcase{!A}{p}
  {$\pSat{v(\Gamma^*)}{p}$ ancak ve ancak
    $\pAssign v(\Gamma^*)(p)=\True$ ise (sağlama tanımı gereği); bu da
    ancak ve ancak $p\in\Gamma^*$ ise geçerlidir
    ($\pAssign v(\Gamma^*)$ kuruluşu gereği).}}

\tagitem{prvNot}{%
  \indcase{!A}{\lnot !B}
    {$\iftag{FOL}{\Sat{M(\Gamma^*)}}{\pSat{v(\Gamma^*)}}{\indfrm}$
    ancak ve ancak
    $\iftag{FOL}{\Sat/{M(\Gamma^*)}{!B}}{\pSat/{v(\Gamma^*)}{!B}}$
    ise (sağlama tanımı gereği). Tümevarım varsayımıyla,
    $\iftag{FOL}{\Sat/{M(\Gamma^*)}{!B}}{\pSat/{v(\Gamma^*)}{!B}}$
    ancak ve ancak $!B\notin\Gamma^*$ ise geçerlidir. $\Gamma^*$
    tutarlı ve !!{complete} olduğundan $!B\notin\Gamma^*$ ancak ve
    ancak $\lnot !B\in\Gamma^*$ ise.}}{}

\tagitem{prvAnd}{%
\iftag{probAnd}{%
  \indcase!{!A}{!B \land !C}{}}{%
  \indcase{!A}{!B \land !C}
    {$\iftag{FOL}{\Sat{M(\Gamma^*)}}{\pSat{v(\Gamma^*)}}{\indfrm}$
    ancak ve ancak hem
    $\iftag{FOL}{\Sat{M(\Gamma^*)}}{\pSat{v(\Gamma^*)}}{!B}$ hem de
    $\iftag{FOL}{\Sat{M(\Gamma^*)}}{\pSat{v(\Gamma^*)}}{!C}$ ise
    (sağlama tanımı gereği); tümevarım varsayımıyla bu, ancak ve ancak
    hem $!B\in\Gamma^*$ hem de $!C\in\Gamma^*$ ise geçerlidir.
    \olref[ccs]{prop:ccs}\olref[ccs]{prop:ccs-and} uyarınca bu koşul,
    ancak ve ancak $(!B\land !C)\in\Gamma^*$ ise geçerlidir.}}}{}

\tagitem{prvOr}{%
\iftag{probOr}{%
  \indcase!{!A}{!B \lor !C}{}}{%
  \indcase{!A}{!B \lor !C}
    {$\iftag{FOL}{\Sat{M(\Gamma^*)}}{\pSat{v(\Gamma^*)}}{\indfrm}$
    ancak ve ancak
    $\iftag{FOL}{\Sat{M(\Gamma^*)}}{\pSat{v(\Gamma^*)}}{!B}$ ya da
    $\iftag{FOL}{\Sat{M(\Gamma^*)}}{\pSat{v(\Gamma^*)}}{!C}$ ise
    (sağlama tanımı gereği); tümevarım varsayımıyla bu, ancak ve ancak
    $!B\in\Gamma^*$ ya da $!C\in\Gamma^*$ ise geçerlidir.
    \olref[ccs]{prop:ccs}\olref[ccs]{prop:ccs-or} uyarınca bu da ancak
    ve ancak $(!B\lor !C)\in\Gamma^*$ ise geçerlidir.}}}{}

\tagitem{prvIf}{%
\iftag{probIf}{%
  \indcase!{!A}{!B \lif !C}{}}{%
  \indcase{!A}{!B \lif !C}
    {$\iftag{FOL}{\Sat{M(\Gamma^*)}}{\pSat{v(\Gamma^*)}}{\indfrm}$
    ancak ve ancak
    $\iftag{FOL}{\Sat/{M(\Gamma^*)}}{\pSat/{v(\Gamma^*)}}{!B}$ ya da
    $\iftag{FOL}{\Sat{M(\Gamma^*)}}{\pSat{v(\Gamma^*)}}{!C}$ ise
    (sağlama tanımı gereği); tümevarım varsayımıyla bu, ancak ve ancak
    $!B\notin\Gamma^*$ ya da $!C\in\Gamma^*$ ise geçerlidir.
    \olref[ccs]{prop:ccs}\olref[ccs]{prop:ccs-if} uyarınca bu da ancak
    ve ancak $(!B\lif !C)\in\Gamma^*$ ise geçerlidir.}}}{}

\iftag{FOL}{%
\tagitem{prvAll}{%
  \iftag{probAll}{%
    \indcase!{!A}{\lforall[x][!B(x)]}{}}{%
    \indcase{!A}{\lforall[x][!B(x)]}
      {$\Sat{M(\Gamma^*)}{\indfrm}$ ancak ve ancak bütün~$t$ terimleri
      için $\Sat{M(\Gamma^*)}{!B(t)}$ ise
      (\olref{prop:quant-termmodel}). Tümevarım varsayımına göre bu,
      ancak ve ancak bütün~$t$ terimleri için $!B(t)\in\Gamma^*$ ise;
      \olref[hen]{prop:saturated-instances} uyarınca bu da ancak ve
      ancak $\lforall[x][!B(x)]\in\Gamma^*$ ise geçerlidir.}}}{}

\tagitem{prvEx}{%
  \iftag{probEx}{%
    \indcase!{!A}{\lexists[x][!B(x)]}{}}{%
    \indcase{!A}{\lexists[x][!B(x)]}
      {$\Sat{M(\Gamma^*)}{\indfrm}$ ancak ve ancak en az bir~$t$ terimi
      için $\Sat{M(\Gamma^*)}{!B(t)}$ ise
      (\olref{prop:quant-termmodel}). Tümevarım varsayımına göre bu,
      ancak ve ancak en az bir~$t$ terimi için $!B(t)\in\Gamma^*$ ise;
      \olref[hen]{prop:saturated-instances} uyarınca bu da ancak ve
      ancak $\lexists[x][!B(x)]\in\Gamma^*$ ise geçerlidir.}}}{}
}{}
\end{enumerate}
\end{proof}

\tagprob[FOL]{probOr,probAnd,probIf,probEx,probAll}

\begin{prob}
\olref[fol][com][mod]{lem:truth} lemmasının ispatını tamamlayın.
\end{prob}

\tagendprob

\tagprob[notFOL]{probOr,probAnd,probIf}

\begin{prob}
\olref[pl][com][mod]{lem:truth} lemmasının ispatını tamamlayın.
\end{prob}

\tagendprob

